\section{Introduction}
\label{sec:introduction}

\subsection*{Context}
In a globalized world defined by constant cross-cultural interaction, effective techniques for Second Language Acquisition (SLA) are essential. This paper introduces a novel educational and diagnostic tool designed for linguistics researchers and students to analyze the structural and cognitive processes involved in learning a new language.

\subsection*{Research Motivation}
The core research question investigates how to simulate SLA using Large Language Models (LLMs) to determine if Transformer-based architectures mirror human learning mechanisms. By tracking a model's trajectory across various training stages, we observe how it bridges semantic gaps, handles grammatical restructuring, and exhibits human-like transitional phases—such as vocabulary mixing—when exposed to a completely foreign linguistic framework.

\subsection*{Project Objectives}
To answer the overarching research question, this project pursues three core objectives:
\begin{itemize}
    \item \textbf{Cross-Lingual Implementation:} Fine-tuning monolingual English GPT-2 models to acquire Spanish while actively mitigating the catastrophic forgetting of their native English syntax and foundational knowledge.
    \item \textbf{Interactive Educational Visualization:} Developing a physics-based gamified medium using the Godot engine to visualize the model's linguistic capacity. The simulation places the user and the AI on a floating boat. An external judge (Gemini 2.5 Flash) evaluates the AI's responses in real-time; poor linguistic scores physically add weight to the boat, challenging the user to adapt their communication to the model's epoch-specific capabilities to prevent sinking.
    \item \textbf{Rigorous Benchmarking:} Implementing a comprehensive benchmarking pipeline to quantitatively test and compare the models' grammatical, syntactical, and semantic evolution across multiple epochs.
\end{itemize}